\documentclass[letterpaper]{article}

\usepackage[letterpaper,margin=1in]{geometry}

% Compile from the repository root so that this path resolves to
% the pinned WSP submodule.
\input{wsp/testing/test-case-library.tex}

\begin{document}

\def\TCID{%
  TC-0002
}

\def\TCTitle{%
  Range rules of ProgressBar and Slider
}

\def\TCPurpose{%
  Verify the rules for Min, Max and Value of a ProgressBar and a Slider: Value
  must lie inside the range (error 380 otherwise), the range must stay
  ordered, a narrower range brings Value inside, and a move by the user stops
  at the ends.
}

\def\TCPriority{%
  \TCPriorityHigh
}

\def\TCRequirementRef{%
  REQ-0010, REQ-0025, REQ-0037
}

\def\TCDesignRef{%
  Not applicable; the rules are in the named header.
}

\def\TCFeatureRef{%
  \texttt{include/wfc/mscomctl/range.hpp} (\texttt{RangeModel}) and
  \texttt{tests/range\_test.cpp} and the runner
  \texttt{tests/tc-0002-range.ps1}
}

\def\TCTestTechnique{%
  Equivalence partitioning and boundary value analysis of the rule functions
  (ISO/IEC/IEEE 29119-4 specification-based techniques).
}

\def\TCPreconditions{%
  The library is built standalone (\texttt{cmake -S . -B out/build}, then
  \texttt{cmake --build out/build --config Debug}) with
  \texttt{WFC\_BUILD\_TESTS} on; CTest runs the test as
  \texttt{TC-0002-range}
  (or \texttt{wfc-mscomctl-TC-0002-range} when WFC adds
  the library).
}

\def\TCEnvironment{%
  Any target the library builds for (Windows x86, x64 and ARM64 with MSVC,
  Linux through Winelib); local runs use the x64 Debug configuration of the
  standalone build. The test needs no window, registry or network.
}

\def\TCAssumptions{%
  The rule header under test is the one built into the library; the test is
  compiled from the same source revision.
}

\def\TCInputData{%
  A model of range 0 to 100 and value 0, and a Slider model of range 0 to 10
  and value 5; values 40, 100.5, -0.5, 30, 50, -1, 500, 0, 99.
}

\def\TCInitialState{%
  A fresh process; the test keeps no state between runs and holds no files or
  handles.
}

\def\TCProcedure{%
  \begin{enumerate}
    \item Run \texttt{mscomctl\_range\_test} (CTest \texttt{TC-0002-range}).
    \item Set Value to 40, again to 40, to 100.5 and to -0.5; set Max to 30;
      set Min to 50, Max to -1, Min to 30 and Max to 0; set Min to 10.
    \item Move to 500 and to 0.
    \item On the Slider model set Value to 99 and to -1.
  \end{enumerate}
}

\def\TCExpectedResult{%
  \begin{enumerate}
    \item Value 40 is valid and changes; the same value again is valid and no
      change; 100.5 and -0.5 are invalid (error 380) and leave Value at 40.
    \item A narrower Max of 30 is valid and brings Value to 30; a Min above
      Max, a Max below Min and a Min or Max equal to the other are invalid and
      leave the range alone; a wider range leaves Value.
    \item A move to 500 stops at Max (no change at 30) and a move to 0 stops
      at Min (Value 10).
    \item A Slider brings a Value of 99 or -1 inside Max or Min instead of
      refusing it.
    \item The program prints \texttt{range model: ok} and exits with code 0.
  \end{enumerate}
}

\def\TCPostConditions{%
  The process ends; nothing is written outside CTest's own log.
}

\TCTemplate

\end{document}
